% Chapter 2: Review of Related Literature

\chapter{Review of Related Literature}
\noindent{This systematic review was conducted following the \ac{PRISMA} 2020 guidelines to identify and synthesize relevant literature for developing a parallelized agent-based optimization framework for jeepney route system planning with integrated passenger behavior simulation. The review identified 71 studies from an initial pool of 162 records spanning January 2021 to November 2025.

% Figure 1: PRISMA Flow Chart
\input{chap2/figures/figure_prisma}

% Table 1: Database Queries
\input{chap2/tables/table_1}

% Table 2: Exclusion Criteria
\input{chap2/tables/table_2}

The systematic literature search was conducted using Google Scholar with three comprehensive database queries as listed in Table 1. Beyond the queried literature, the search for supplementary records for the literature review yielded 6 additional records. After removing duplicates, 156 records were found to be eligible for screening. The three researchers screened these records, and excluded 97 based on the title and abstract. The criteria for exclusion and sample citations of each are detailed in Table 2. After performing a final full-text assessment of these articles, a total of 71 papers were chosen to be included in this systematic review.}

% Chapter 2 Section 2.1: The Philippine Public Transport Landscape: Challenges and Operational Context

\section{The Philippine Public Transport Landscape: Challenges and Operational Context}
\noindent{This section establishes the problem space of Philippine public transport: a road network dominated by fragmented paratransit such as jeepneys, tricycles, e-trikes, operating under severe congestion, operational inefficiency, and mounting environmental pressures. The literature demonstrates that intersections routinely fail to handle mixed traffic flows, first- and last-mile services operate without coordination, and small \ac{PUV} are indispensable yet deeply inefficient. These conditions justify the need for data-driven optimization of paratransit routes and operations tailored to local conditions.}

% Chapter 2 Section 2.1 Subsection 2.1.1: Operational Inefficiencies & Traffic Management

\subsection{Operational Inefficiencies and Traffic Management}
\noindent{\citet{lim2021} documented the Banlic-Mamatid intersection in Cabuyao City operating at Level of Service F, the worst traffic operating conditions with delays averaging 104 seconds per vehicle, driven by signal misalignment and narrow approaches that create cascading bottlenecks. This reflects a broader national pattern where intersections designed for simpler traffic now struggle with complex, mixed-mode flows. \citet{lim2022}  identify the first and last mile, which refers to the short trips linking homes to major transport routes and final destinations, as a critical weakness. Uncoordinated tricycles, jeepneys, and e-trikes operate without integration to formal systems. Understanding these delays matters because \citet{sobisol2025} used system dynamics to show how terminal dwell times, or the time vehicles spend waiting or loading at terminals, road congestion, and irregular dispatching compound cumulatively, underscoring that time-based metrics are essential to evaluating operations.

Tricycles and e-trikes function as de facto feeder services that connect passengers to main transport routes, although they operate with suboptimal efficiency. \citet{nacion2022} found that e-trikes have lower operating costs but face range constraints, while conventional tricycles exhibit high fuel consumption and idle times despite widespread deployment. \citet{salison2025} estimated that the tricycle sector alone generates significant energy demand and emissions in cities like Quezon City. Localized traffic simulation tools have proven necessary for testing solutions. \citet{eden2021} validated LocalSim, a simulator tuned to Philippine driving behavior, demonstrating that locally calibrated parameters improve realism compared to generic models. \citet{eugenio2023} used PTV VISSIM to evaluate infrastructure improvements, while \citet{lim2021} tested signal optimization strategies, all revealing that while operational tweaks help, geometric constraints often limit improvement; network-level redesign is needed.}

% Chapter 2 Section 2.1 Subsection 2.1.2: Socio-Economic & Environmental Drivers

\subsection{Socio-Economic and Environmental Drivers}
\noindent{\citet{vergel2022} estimated that road transport dominates national energy use, with \ac{PUV} consuming substantial shares and demand projected to rise. \citet{salison2025} confirm tricycles are significant local emitters at scale. \citet{rizki2020} provide methodology for calculating urban transport emissions, enabling translation of operational changes into emission outcomes. Economic pressures compound these challenges: \citet{tuano2021} showed that fuel excise taxes increase poverty incidence. Simultaneously, \citet{roquel2022} demonstrated that fuel price spikes shift traveler behavior and threaten operator viability. For small operators with thin margins, inefficiency becomes a financial crisis during fuel price shocks.

External shocks expose system vulnerability. \citet{ancheta2023} found that \ac{COVID-19} disproportionately disrupted low-income commuters in Metro Manila, forcing mode shifts and reliance on poorly coordinated paratransit. \citet{chen2025} showed that telecommuting benefits affluent workers while essential workers cannot reduce travel, raising equity concerns. These findings underscore that Philippine paratransit must be resilient and equitable in that it is adaptable to demand shifts and external shocks while protecting vulnerable users. The convergence of these operational, economic, and social pressures establishes the imperative for integrated, data-driven optimization grounded in local conditions and evaluated across time, energy, emissions, and equity dimensions.}

% Chapter 2 Section 2.2: Passenger Behavioral Modeling: The Human Element

\section{Passenger Behavioral Modeling: The Human Element}
\noindent{This section defines Behavioral Simulation by establishing what passenger behavior looks like in Philippine paratransit systems. Passengers make decisions based on walking tolerances, waiting preferences, safety concerns, and satisfaction thresholds. Understanding these behavioral parameters is essential for creating realistic agents that respond to route design changes. The literature reveals that passenger satisfaction depends on multiple factors such as travel time, accessibility, transfer convenience, and perceived safety. These factors must be integrated into the simulation’s fitness function to reflect how commuters respond to alternative paratransit configurations.}

% Chapter 2 Section 2.2 Subsection 2.2.1 : Satisfaction and Service Quality

\subsection{Satisfaction and Service Quality}
\noindent{Passenger happiness extends beyond faster commute times alone. \citet{jou2023} examined commuter satisfaction with public utility buses in Occidental Mindoro, revealing that satisfaction stems from reliability, comfort, safety, and responsiveness rather than speed alone. \citet{mallillin2020} analyzed Philippine commuter behavior, showing that passengers make mode choices based on broader service quality perceptions, establishing that fitness function parameters must capture dimensions beyond travel time. For informal paratransit like jeepneys, which refers to flexible and semi-regulated transport services such as jeepneys, \citet{romero-torres2023} measured service quality in Mexican mototaxis and found that passengers value frequency, safety, and courtesy, demonstrating that Quality of Service in informal paratransit requires multidimensional assessment.}

% Chapter 2 Section 2.2 Subsection 2.2.2 : Walking, Accessibility, and Inclusivity

\subsection{Walking, Accessibility, and Inclusivity}
\noindent{\citet{ferels2025} developed a dynamic walkability model showing how walking distances, topography, and safety influence stop access. Critically, accessibility varies across groups. \citet{gazarin2024} and \citet{mcginnis2025} showed that gender, safety concerns, and household constraints create distinct mobility patterns, establishing that simulation must include diverse agent types, not one generic agent but populations reflecting gender, age, and mobility variation. \citet{macfarlane2021} estimated activity patterns for wheelchair users, revealing accessibility constraints for vulnerable groups. \citet{agramon2024} demonstrated that passenger safety concerns influence route choice, particularly for vulnerable populations.}

% Chapter 2 Section 2.2 Subsection 2.2.3 : Transfer Aversion: Wait vs. Detour

\subsection{Transfer Aversion: Wait vs. Detour}
\noindent{\citet{lu2024} mathematically formalized the trade-off between waiting time and detours, establishing the basis for modeling transfer aversion in agents, which refers to passengers’ tendency to avoid switching between vehicles or transport modes. \citet{malichova2020} conducted Value of Travel Time analysis, a method used to measure how passengers value time spent traveling, waiting, or transferring, revealing that passengers exhibit heterogeneous time valuations in which some tolerate longer journeys to avoid transfers, while others accept transfers to minimize total time. These parameters must embed realistic trade-off behavior into agent decision-making.}

% Chapter 2 Section 2.2 Subsection 2.2.4 : Modern Urban Structure & Agent Flow

\subsection{Modern Urban Structure and Agent Flow}
\noindent{\citet{zhangb2024} analyzed monocentric city structures and transit layouts, justifying careful city structure specification in simulation parameterization. \citep{david2022transport} further examined transport policies in polycentric cities, showing that multiple urban centers create more distributed travel demand patterns and require coordinated transport planning across interconnected districts. \citep{roth2011structure} analyzed urban movement structures and demonstrated that polycentric activity patterns generate entangled hierarchical flows between different activity centers, reinforcing the importance of modeling realistic movement interactions within cities. \citet{burke2025} presented a 3D agent-based model for urban densification, providing a framework for how agents interact with high-density centers including congestion dynamics. \citet{sonnenschein2022} provided guidance on linking spatial environments (\ac{GIS}) to agent behavior, establishing methodologies for grounding agents in realistic geography, street networks, and land-use patterns. \citep{fielbaum2016optimal} further showed that optimal public transport networks depend heavily on urban structure, where both monocentric and polycentric forms influence route efficiency, connectivity, and passenger movement behavior.}

% Chapter 2 Section 2.3 : Modernizing Paratransit: From Informal to Shared Mobility

\section{Modernizing Paratransit: From Informal to Shared Mobility}
\noindent{This section bridges the gap between current jeepney operations and modern shared mobility solutions. It reviews how the rest of the world handles Shared Mobility and \ac{DRT}, framing the study as modernization of the jeepney. Critically, this section addresses a key reviewer concern by emphasizing that operator adoption relies on economic incentives, policy support, and demonstrated financial stability.}

% Chapter 2 Section 2.3 Subsection 2.3.1 : The Transition to Shared / Smart Mobility

\subsection{The Transition to Shared / Smart Mobility}
\noindent{Jeepneys can evolve into Shared Mobility and smart transport systems by adopting digital coordination and real-time routing. \citet{seng2023} describe how ridesharing and crowdsourcing platforms use digital matching and smartphone interfaces to pool trips and reduce empty running, positioning these as core technologies for urban transport rather than alternatives to public transit. This demonstrates a transition from traditional public transport operations toward more coordinated shared mobility systems. \citet{barlis2022} applies this to Philippine contexts, using hybrid agent-based and system dynamics modeling to analyze how policies such as subsidies, charging infrastructure, and regulatory support influence the adoption of electric jeepneys and their transition toward coordinated fleets. These sources show that modernization is technically feasible and can improve efficiency if paired with the right policy environment.

However, modernization in developing countries follows different patterns than in the West. \citet{zeeshan2022} review bus network design in developing nations and highlight that Western planning models assume formal institutions and strict compliance, which often do not match the flexible, informal reality of Global South transit systems. \citet{zeeshan2023}  proposes a reliability-based priority system that accounts for operator incentives and elastic demand, arguing that realistic design must work within local constraints rather than impose idealized structures. Directly addressing the driver adoption question, \citet{saxena2023} examined why rickshaw operators in Bhopal, India, willingly adopt battery-operated vehicles, finding that adoption depends on perceived profitability, upfront costs, and policy credibility. Together, these studies suggest that jeepney modernization toward shared and smart mobility will succeed only if new systems protect or improve driver income and are backed by stable, trustworthy policies that operators find believable.}

% Chapter 2 Section 2.3 Subsection 2.3.2 : Demand Responsive Transport (DRT) and Ride Pooling

\subsection{Demand Responsive Transport (DRT) and Ride Pooling}
\noindent{Beyond vehicle technology, route structures themselves must evolve toward flexibility and demand responsiveness. \citet{pavanini2023} and \citet{pavanini2025} present case studies of rural and inner-area \ac{DRT} in Italy, showing that flexible, on-demand routing can sustain service in low-density areas where fixed routes would be inefficient and financially unviable, by adjusting vehicle paths to actual passenger requests rather than adhering to rigid schedules. This model is directly applicable to Philippine last-mile and dispersed urban corridors where traditional jeepney routes often run with low occupancy. \citet{zwick2022} reviews ride-pooling implementations, a system in which multiple passengers share a vehicle traveling along similar routes. They concluded that pooling significantly increases vehicle occupancy and reduces vehicle-kilometers traveled, which refers to the total distance traveled by vehicles within a transport network. \citet{shen2024} demonstrates that integrating mobility-on-demand, a transport service where vehicles are requested as needed by passengers, with traditional public transport in suburban areas proves more cost-effective. It is more cost-effective than relying solely on fixed routes. These findings establish that dynamic routing is operationally superior to static routes in many contexts.

Making \ac{DRT} practical requires robust operational strategies. \citet{kim2020} evaluates \ac{DRT} systems using optimal routing algorithms that balance waiting time, detour distance, and operator efficiency, providing a clear framework for how systems can serve passengers without unsustainable costs. \citet{gokay2021} demonstrates that using designated meeting points in real-time ride sharing, where passengers are grouped dynamically during active operations, keeps systems scalable while maintaining reasonable walking distances for passengers. Combined with evidence on operator adoption, these studies paint a coherent vision: modernized jeepney systems should be demand-responsive and pooled, where routes and stops adjust to demand patterns while remaining financially viable for drivers, directly aligning modernization with the simulation-based optimization this study proposes.}

% Chapter 2 Section 2.4 : Agent-Based Modeling (ABM) in Transportation

\section{Agent-Based Modeling in Transportation}
\noindent{\ac{ABM} provides a computational framework for jeepney route optimization that fundamentally differs from traditional static network design. Conventional approaches assume equilibrium conditions and homogeneous commuter behavior. In contrast, \ac{ABM} simulates heterogeneous agents, including individual commuters and vehicle operators, whose localized interactions generate emergent routing patterns that reflect real operational dynamics. This capability directly addresses the gap between theoretical route optimization and commuter-centered objectives. It enables simultaneous modeling of individual behavioral responses such as transfer aversion, boarding delays, and wait time sensitivity. It also captures system-level feedback including congestion propagation and demand shifts. For Philippine paratransit systems characterized by informal operations and high demand elasticity, \ac{ABM} proves particularly valuable in capturing how routes emerge from passenger choice dynamics rather than centralized planning alone.}

% Chapter 2 Section 2.4 Subsection 2.4.1 : ABM in Data-Scarce Environments

\subsection{Agent-Based Modeling in Data-Scarce Environments}
\noindent{Integrating agent-based microsimulation with activity-based demand modeling provides a methodological foundation for capturing both behavioral responses and spatiotemporal network feedback. \citet{lim2023} establish this integration framework, demonstrating that activity-based models, which estimate travel demand based on people’s daily activities and travel behavior, alone cannot represent how passengers respond to changing network conditions. The limitation is critical for optimizing demand-responsive services. Agent-based microsimulation overcomes this by implementing autonomous agents that interact dynamically within networks, enabling spatiotemporal optimization of supply and demand. Yet \ac{ABM} applications have historically required extensive travel survey data. \citet{elgohary2025} demonstrate feasibility in data-constrained contexts by constructing \ac{ABM} using social media platforms to extract behavioral travel patterns, generating synthetic agent populations without costly surveys. This approach is directly applicable to Philippine cities where \ac{OD} surveys, which track where trips begin and end, are prohibitively expensive but alternative data traces exist. \citet{manser2020} extend this to large-scale networks in Zurich, validating that \ac{ABM} incorporating dynamic demand response, where transport demand changes according to network conditions and passenger behavior, produces operationally superior solutions with fewer lines, higher frequencies, smaller vehicles, yet higher predicted ridership. Their findings establish that sophisticated route optimization is computationally feasible and yields outcomes aligned with both service efficiency and commuter experience.}

% Chapter 2 Section 2.4 Subsection 2.4.2 : Simulation Dynamics

\subsection{Simulation Dynamics}
\noindent{Emergent system behavior from agent interactions cannot be predicted through aggregated or separate modeling approaches. \citet{chandiramani2025} identify contemporary \ac{ABM} applications for public transportation including adaptive vehicle routing and real-time signal integration, establishing operational viability standards. \citet{pan2025} quantify impacts of multi-agent traffic optimization on travel behavior, demonstrating that agent-level modifications propagate to measurable system-level efficiency gains. \citet{torre2021} validate this in Metro Manila’s rapid transit system, finding that simultaneous modeling of passenger crowds and operations generates dynamics impossible to predict from separate models; validation against smart card and video data confirms behavioral fidelity and reveals that operational parameters such as train frequency below threshold produce substantial increases in passenger travel time. These validation methodologies, combining transaction data, video observation, and commuter surveys, directly inform how jeepney \ac{ABM} validates against boarding-alighting observations and \ac{GPS} traces. \citet{koch2020} demonstrate \ac{ABM} effectiveness in evaluating network vulnerability under adverse conditions. \citet{dytckov2020} integrate classical four-stage demand modeling with agent-based simulation for responsive transport, disaggregating aggregate data into individual trips with 3-–20\% validation accuracy against commuting records, directly applicable to jeepney demand where aggregated boarding counts must be translated into realistic passenger-level \ac{OD} patterns. \citet{liu2025} propose \ac{LLM}-agent frameworks incorporating machine learning for enhanced behavioral realism. They position current behavioral simulation approaches including commute-time-minimizing pathfinding with transfer penalties and boarding costs as the contemporary standard before \ac{AI} fundamentally transforms transportation modeling.}

% Chapter 2 Section 2.5 : Optimization Algorithms and Network Design 

\section{Optimization Algorithms and Network Design}
\noindent{Agent-based optimization of jeepney routes requires a computational solver capable of searching large, combinatorial route spaces while evaluating candidate solutions under realistic passenger demand and behavioral constraints. \citet{nnene2023} establish simulation-based optimization as a standard paradigm for transit network design, demonstrating that metaheuristics coupled with agent-based simulators form an effective ``optimize $\rightarrow$ simulate $\rightarrow$ evaluate'' loop necessary when stochastic demand and detailed passenger interactions must be represented. This study’s \ac{GA}+\ac{ACO} framework directly instantiates this paradigm: genetic and ant-inspired operators generate jeepney route variants, the \ac{ABM} simulates passenger trips under those routes, and Total User Cost (incorporating travel, waiting, and transfer penalties) feeds back to guide the search. The approach is thus grounded in contemporary best practice for transit optimization rather than relying on closed-form formulations that neglect passenger-centric metrics and informal operational dynamics characteristic of Philippine paratransit.}

% Chapter 2 Section 2.5 Subsection 2.5.1 : Network Design Problems 

\subsection{Network Design Problems}
\noindent{\citet{guillermo2022a} and \citet{guillermo2022b} demonstrate that strategic transit route design can be operationalized through logit-based scoring of candidate routes using travel time, waiting time, and transfer penalties, while representing routes, stops, and \ac{OD} relations as graph-structured data enabling flexible querying and evaluation. These two complementary works validate both the scoring paradigm (Total User Cost evaluation via the TravelGraph cost function) and the graph representation used in this study to encode walking, boarding, and jeepney-travel costs. \citet{hatzenbuhler2022} and \citet{borchers2024} reinforce that network design in modern transit optimization balances passenger-oriented objectives (generalized travel time, accessibility) against operational and sustainability criteria, and that hybrid simulation--metaheuristic methods are the current standard for solving such multi-faceted design problems. These works position the \ac{GA}+\ac{ACO} solver as a methodologically sound approach for exploring jeepney network variants while the embedded \ac{ABM} evaluates outcomes from the commuter perspective.}

\subsection{Optimization Strategies (Heuristics and Fleet)}
\noindent{In paratransit systems, once the spatial route trajectories are determined, the secondary task of operational resource distribution becomes crucial. \citet{guillermo2022a} and \citet{guillermo2022b} establish that network performance depends heavily on fleet assignment constraints, showing that even spatially optimized route layouts can fail if passenger capacity and vehicle dispatch frequencies are poorly matched to localized demand. This highlights the importance of integrating a mathematical fleet allocation model within the transit route network design problem (TRNDP) environment. \citet{hatzenbuhler2022} and \citet{borchers2024} reinforce that operational parameters—such as vehicle capacities, frequency settings, and total fleet sizing—must be balanced alongside passenger-oriented objectives like travel time minimization and accessibility. To achieve this balance in decentralized, uncoordinated paratransit systems, this study implements Mohring's square-root allocation rule to distribute a fixed fleet of vehicles dynamically across evolved routes based on passenger demand. This hybrid approach ensures that the genetic algorithm evolves robust spatial route topologies, while the operational fleet allocation maximizes service frequencies along high-demand corridors, thereby directly minimizing passenger waiting times and the overall Total User Cost.}

% Chapter 2 Section 2.6 : Hybrid and Parallel and Metaheuristics

\section{Hybrid and Parallel Metaheuristics}
\noindent{Jeepney route optimization requires combining multiple metaheuristic mechanisms to balance global search efficiency with local solution refinement under complex, combinatorial problem structures. Single-algorithm approaches face trade-offs: \ac{GA} excel at global exploration but may converge slowly without sufficient local feedback, while \ac{ACO} provides strong local search through pheromone-based attraction but risks premature convergence in early iterations. \citet{shi2022} and \citet{ran2024} demonstrate that hybrid \ac{GA}+\ac{ACO} architectures overcome these limitations by leveraging complementary strengths, \ac{GA}’s population diversity and adaptive operators paired with \ac{ACO}’s positive feedback mechanism. These approaches reduced convergence iterations by up to 47\% and improved solution quality compared to single-method baselines. This study adopts \ac{GA}+\ac{ACO} specifically to mutate jeepney route sets using \ac{GA} while using pheromone-influenced path construction through \ac{ACO} to refine candidate solutions, all evaluated in parallel via \ac{ABM} simulation to enable rapid exploration of the route configuration space.

Recent studies further support the integration of surrogate-assisted evaluation and Lamarckian learning within memetic optimization for combinatorial routing and scheduling problems. \citep{behmanesh2021comparison} demonstrated that memetic algorithms combining evolutionary search with embedded local refinement outperform conventional \ac{GA} in integrated logistics network optimization, validating the effectiveness of inheritance-based local improvement strategies in complex routing structures. Similarly, \citep{liang2021surrogate} showed that surrogate-assisted optimization significantly reduces computational burden in large-scale traffic network optimization by replacing expensive simulation evaluations with predictive surrogate models while maintaining solution quality. These findings support the proposed Surrogate-Assisted Memetic Algorithm with Lamarckian learning by demonstrating that learned local improvements can be inherited across generations and that surrogate-assisted evaluation can accelerate optimization in computationally expensive transportation systems.}

% Chapter 2 Section 2.6 Subsection 2.6.1 : Why Hybrid (GA + ACO)

\subsection{Why Hybrid (GA + Ant Colony Optimization)}
\noindent{\citet{shi2022} provide empirical evidence that improved hybrid \ac{GA}+\ac{ACO} reduces iterations by 46–47\% and turns by 21–75\% on grid-based path planning compared to either algorithm alone, showing that genetic crossover and mutation complement ant pheromone updates in balancing exploration and exploitation. \citet{ran2024} extend this paradigm to fuzzy-enhanced hybrid \ac{GA}+\ac{ACO} for automated parameter tuning, demonstrating that adaptive fuzzy control of \ac{ACO} parameters and \ac{GA} crossover operators further improves convergence while avoiding the brittleness of hand-tuned thresholds. These validations, spanning path planning and clustering optimization, establish that \ac{GA}+\ac{ACO} hybrids remain state-of-the-art for combinatorial route design, positioning this study’s adoption of the hybrid for jeepney route mutation and evaluation as methodologically current and empirically justified for minimizing both solution time and the Total User Cost across multiple urban configurations.}

\subsection{Non-Standard Pheromone Mechanisms in Hybrid Frameworks}
\noindent{Canonical \ac{ACO} constructs complete solutions iteratively, relying exclusively on positive reinforcement. However, modern routing problems require complex evolutionary hybrids that utilize non-standard pheromone mechanisms. \citep{nurcahyadi2022negative} established a formal algorithmic framework for negative learning in \ac{ACO}, demonstrating that increasing negative pheromone values on sub-optimal path components prevents algorithmic stagnation. This repulsion mechanism forces the localized exploration of unrouted corridors that standard \ac{ACO} would ignore.

Furthermore, integrating \ac{ACO} within a broader evolutionary algorithm alters its functional role. Rather than constructing entirely new routes, pheromone-guided searches operate highly effectively as localized mutation operators. In the smartPATH framework, a hybrid \ac{GA}-\ac{ACO} approach utilizes genetic operators for global search while restricting \ac{ACO} to localized path refinement, demonstrating superiority in convergence speed over isolated algorithms \citep{chari2012smartpath}. Similarly, \citep{zhao2026traversal} embedded genetic mutation operators directly into an \ac{ACO} framework to dynamically sustain population diversity during path planning.

To permanently capture these localized improvements, hybrid frameworks employ specific inheritance strategies. Memetic computing establishes that local search improvements must be integrated back into the global search population. Meta-Lamarckian learning structures coordinate this by directly overwriting the global chromosome with the localized heuristic solution \citep{neri2012meta}. By applying \ac{ACO} strictly as a localized heuristic evaluator whose refined paths undergo Lamarckian inheritance back into the global \ac{GA}, the framework maximizes both global exploration diversity and local exploitation precision.}

% Chapter 2 Section 2.7 : Comparative Technologies and Data-Driven Approaches

\section{Comparative Technologies and Data-Driven Approaches}
\noindent{This study deliberately employs agent-based simulation coupled with metaheuristic route optimization rather than pure machine learning or autonomous fleet architectures, reflecting the distinct requirements of informal transit systems in developing contexts. Machine learning approaches, including mode choice prediction via XGBoost and neural networks \citep{tran2020, ma2023} and \ac{GPS}-based behavioral analysis \citep{joseph2020}, excel at inferring travel patterns and demand elasticity. However, these methods operate within fixed network structures rather than generating alternative route topologies. Meinhardt and Nagel (2025) extend this to dynamic pricing models that optimize demand response, but these models also focus on within-network optimization  rather than structural design. Although these predictive methods are valuable for understanding commuter behavior, they cannot systematically explore which route configurations minimize Total User Cost. Autonomous fleet and robotaxi systems also require extensive real-time data, advanced infrastructure, and highly formalized transport operations that are not yet feasible in Philippine paratransit systems. Therefore, agent-based simulation and metaheuristic optimization remain more appropriate for resource-constrained environments where large-scale data collection is limited. }

% Chapter 2 Section 2.7 Subsection 2.7.1 : Machine Learning and Prediction

\subsection{Machine Learning and Prediction}
\noindent{\citet{tran2020} and \citet{ma2023} demonstrate that machine learning, particularly ensemble methods and causal inference, can effectively predict mode choice and estimate unbiased policy effects, validating the use of behavioral data for understanding commuter preferences, though their focus remains predictive rather than generative for new network designs. \citet{joseph2020} shows that \ac{GPS} trajectory data captures realistic movement patterns and reveals behavioral heterogeneity across commuters, aligning with how this study’s survey-validated behavioral parameters such as boarding cost, transfer penalty, walking cost are grounded in actual traveler decision-making. \citet{meinhardt2025} applies transport models to price optimization, demonstrating the power of data-driven control, but again within fixed infrastructure. In contrast, this study requires generating and evaluating candidate route structures that do not yet exist, necessitating simulation-driven exploration rather than within-network prediction alone.}

% Chapter 2 Section 2.7 Subsection 2.7.2 : Autonomous & Future Systems

\subsection{Autonomous and Future Systems}
\noindent{The global research trajectory emphasizes autonomous vehicles, robotaxis, and centrally-controlled shared autonomous vehicle fleets \citep{carreyre2023, tiwari2024, paithankar2024b, sirikhan2022, shobha2023, cheng2024, rath2022}. These technology-centric approaches assume mature data infrastructure, vehicle homogeneity, and capital-intensive electrification that are not yet present in Philippine paratransit. In contrast, this study adopts a human-centric optimization philosophy. Rather than deferring improvement to autonomous maturity, it systematically optimizes existing, manually-operated jeepney fleets through \ac{GA}+\ac{ACO}-driven route exploration and simulation-validated evaluation. This makes the approach directly applicable to the operational reality of informal transit in developing cities and extensible to automation as infrastructure and data maturity increase.}

% Chapter 2 Section 2.8 : Synthesis and Research Gap

\section{Synthesis and Research Gap}
\noindent{Recent literature establishes three foundational validations for this study’s approach, yet reveals a critical integration gap that justifies the novel framework. \citet{zhangb2024} confirm that monocentric urban structures, characterized by concentrated employment and residential zones with directional demand peaks are persistent geographic realities in developing cities, validating the spatial demand patterns that this study models. \citep{david2022transport}, together with \citep{roth2011structure}, further demonstrate that even polycentric urban systems exhibit structured movement hierarchies and interconnected activity flows, reinforcing the importance of accurately modeling urban spatial organization in transport optimization. \citet{shi2022} demonstrate empirically that hybrid \ac{GA}+\ac{ACO} reduces convergence iterations by 46–47\% and improves solution quality compared to single-algorithm approaches, establishing metaheuristic validity for combinatorial route optimization. \citet{elgohary2025} show that agent-based models can be constructed and validated in data-scarce contexts using alternative data sources such as social media, mobile traces, removing a traditional barrier to \ac{ABM} deployment in developing nations lacking comprehensive travel surveys. However, no existing study integrates these three validated elements, namely monocentric demand structure, hybrid \ac{GA}+\ac{ACO} efficiency, and data-scarce \ac{ABM} capability, into a parallel optimization framework specifically designed to redesign fragmented paratransit networks. Furthermore, \citet{lu2024} quantify transfer aversion as a measurable behavioral parameter that substantially reduces transit utility in multimodal journeys, yet prior jeepney optimization work has largely neglected this constraint in route design. This study closes this gap by coupling hybrid \ac{GA}+\ac{ACO} route exploration with \ac{ABM} passenger simulation that explicitly embeds empirical transfer penalties. This enables systematic, user-cost-centric network redesign for informal transit in data-scarce developing contexts. Such capability remains absent from prior literature and directly applicable to Philippine paratransit modernization. While many transport optimizations include travel-time minimization in general, the literature remains underexplored in demonstrating Total User Cost–centric fitness as a system-level objective for informal jeepney network reconstruction under empirically calibrated passenger behaviors and agent-based microsimulation. This limitation is particularly evident in Philippine settings where \ac{ABM}-based evaluation is largely absent. This study therefore positions Total User Cost as a baseline, testable performance objective for informal multi-route redesign, intended to establish whether optimizing this system-level cost yields consistent, measurable improvements under realistic passenger–vehicle interactions.}

\section{Comparative Matrix Analysis}

\FloatBarrier 

\begin{figure}[htbp]
    \centering    
    \includegraphics[width=1\textwidth]{chap2/figures/competitive matrix.jpg}
    \caption{Comparative Matrix of Jeepney Route Optimization Studies}
    \label{fig:compe-matrix}
\end{figure}

\FloatBarrier 

\noindent{Figure~\ref{fig:compe-matrix} presents a comparative matrix of 15 key studies on jeepney and paratransit route optimization, classified by each study’s context, demand modeling approach, behavioral assumptions, simulation framework, optimization algorithm, objectives, and validation method. This comparison reveals clear patterns in the literature. Most studies rely on aggregate \ac{OD} based demand models or other zonal/direct demand estimates rather than simulating individual passenger agents. Behavioral elements like transfer penalties or walking tolerances are usually simplified or fixed, with few works calibrating these factors to observed data. Optimization methods are predominantly single-heuristic: the majority of studies apply one metaheuristic such as \ac{GA} or an \ac{ACO} in isolation, which can limit the balance of global exploration and local exploitation. Indeed, while \ac{GA} excel at broad search and ant-colony methods at path refinement. However, using them alone forces a trade-off. Hybrid approaches can mitigate this limitation. Objectives across the studies tend to focus on travel-time or cost minimization, often a single-objective formulation, and only a few include multi-objective considerations such as balancing efficiency with equity or coverage. Simulation fidelity varies: although simulation-based optimization is increasingly seen as essential for transit network design, only a subset of these studies employ detailed simulation, and even fewer use full agent-based models. Notably, none of the reviewed jeepney route design studies deploy an agent-based microsimulation in a Philippine setting, and many evaluations instead use static assignments or simplified models. Contextually, most prior works address formal or idealized transit scenarios; they often assume a level of institutional structure and compliance that doesn’t reflect the informal, flexible operations of jeepneys. This underscores a gap in applying modern optimization-simulation approaches to the unique context of Philippine paratransit.}

\indent{These observed patterns point to critical gaps in the current literature. First, the scarcity of hybrid metaheuristics is evident. Only rare examples \citep{shi2022} combine algorithms like \ac{GA}+\ac{ACO} for route optimization, despite evidence that such hybrids can improve convergence speed and solution quality over single-algorithm methods. The near-absence of hybrid approaches means that complementary strengths of different heuristics remain underexploited in existing studies. Second, there is a lack of empirically grounded behavioral modeling. Few studies incorporate calibrated traveler behavior parameters. For instance, transfer aversion in which the extra “cost” passengers associate with making a mode transfer is often neglected. This is a significant omission given that transfer penalties can substantially reduce a route’s utility to riders. Without calibrating such factors such as using local survey data, models risk mis-estimating ridership and optimal routes. Third, agent-based simulation is rarely applied, especially in the jeepney context. Even though agent-based models are recognized as powerful tools for capturing individual interactions and stochastic demand, no existing study was found that uses an \ac{ABM} to evaluate jeepney network designs in the Philippines. This gap may stem from data scarcity or computational complexity, but it leaves the person-level dynamics of route changes underexplored. Finally, scalability and validation limitations are evident. Many prior optimizations were tested on small networks or hypothetical scenarios, and few leveraged parallel computing or city-scale \ac{GIS} data. As a result, the literature lacks demonstrations of methods that can handle full city-scale route networks under realistic conditions. In short, no single study integrates advanced behavioral calibration, agent-based simulation, and hybrid optimization into a unified, scalable framework for informal transit. This integration gap has been explicitly noted in recent analyses.}

This study directly addresses these gaps by introducing a novel, integrated approach that combines an empirically calibrated \ac{ABM} with a hybrid metaheuristic in a parallel, city-scale optimization loop. Specifically, the proposed framework couples an agent-based passenger simulation, which models individual commuters’ choices and movements with a hybrid \ac{GA}+\ac{ACO} solver. This integration allows the search process to exploit \ac{GA}’s global exploration and \ac{ACO}’s local path-finding strengths simultaneously, a strategy largely absent from previous jeepney studies. The agent-based model is calibrated using local survey data to incorporate realistic behavior, including actual walking distance tolerances and transfer aversion levels as quantified by recent work. This grounds the optimization in real commuter preferences. Unlike earlier models that assume ideal behavior, the simulation in this thesis penalizes inconvenient transfers and long walks in line with observed traveler discomfort, ensuring that proposed routes remain attractive to passengers. The entire optimize–simulate–evaluate cycle is executed in parallel on multi-core hardware, dramatically accelerating the evaluation of candidate network designs. Such parallelization is known to alleviate computational bottlenecks and makes it feasible to explore large urban networks, which is crucial for city-wide jeepney route optimization. Finally, the study applies this framework to a city-scale, \ac{GIS}-based road network using the actual geography and demand patterns of a Philippine city rather than a toy model, providing a rigorous test of scalability and practical relevance. By integrating agent-based simulation, hybrid \ac{GA}–\ac{ACO} optimization, and parallel computation on a real-world network, this work offers a uniquely comprehensive solution. It moves beyond the limitations of past studies and demonstrates a path forward for data-driven jeepney network redesign, directly supporting the goals of jeepney modernization with a method tailored to the local context and its constraints. The result is a significant step toward filling the identified research gaps and advancing the state of the art in paratransit route optimization.